\section{Primary-source audit: exact statements and boundaries}
This is a targeted audit, not a claim that all relevant literature has
been exhaustively searched. The following locators were checked in the
primary texts or their attached source files during this run.

\paragraph{Kadish--Landsberg.}
\emph{Padded polynomials, their cousins, and geometric complexity
theory}, Communications in Algebra42 (2014),2171--2180,
\href{https://arxiv.org/abs/1204.4693}{arXiv:1204.4693},
Theorems1.2,1.3,1.7 and Proposition1.8.
The first two give padding restrictions and ranges with no ideal copies;
Theorem1.7 describes a generalized Foulkes-map kernel, while
Proposition1.8 describes the relevant normalization under its hypotheses.
Remark1.4 explicitly allows an ideal to contain only some copies of a
Schur module. None of these statements computes the two determinant
evaluation ranks required in the main question. The main article was
read online because the attached KL directory contains definitions and
bibliography but not its main source.

\paragraph{B\"urgisser--Ikenmeyer--Panova and Ikenmeyer--Panova.}
\emph{No occurrence obstructions in geometric complexity theory},
\href{https://arxiv.org/abs/1604.06431}{arXiv:1604.06431v3},
Theorem1.4, excludes occurrence obstructions in its range $n\ge m^{25}$;
section1(c) explicitly retains multiplicity obstructions as a possibility.
Theorem2.1 and Theorem4.9 specify the length and padding restrictions.
The paper's displayed permanent uses internal padding; the separate
variable convention must be handled as in S-004 rather than silently
importing its smaller support cap.
\emph{Rectangular Kronecker coefficients and plethysms in geometric
complexity theory}, Advances in Mathematics319 (2017),40--66,
\href{https://arxiv.org/abs/1512.03798}{arXiv:1512.03798v2},
Theorem1.6 treats ordinary rectangular Kronecker vanishing for
$n>3m^4$; Corollary1.9 gives its degree lower bound. These are not
equalities for actual determinant orbit-closure multiplicities.

\paragraph{D\"orfler--Ikenmeyer--Panova.}
\emph{On geometric complexity theory: multiplicity obstructions are
stronger than occurrence obstructions},
\href{https://arxiv.org/abs/1901.04576}{arXiv:1901.04576v1},
Theorem2.3, compares Chow varieties with bounded-Waring-rank varieties.
Its finite case at degree7 and weight $(34,6,2)$ gives multiplicities
7 and8. Proposition3.2 and section6 contain the explicit finite
calculation; Theorem3.5 supplies its infinite plethysm identity.
These are genuine precedents for full multiplicity comparisons, but
the varieties are different from the padded-permanent/determinant pair.

\paragraph{Ikenmeyer--Kandasamy.}
\emph{Implementing geometric complexity theory: on the separation of
orbit closures via symmetries},
\href{https://arxiv.org/abs/1911.03990}{arXiv:1911.03990v1},
Proposition4.1, Theorems4.2 and4.3, and section5.
Theorem4.2 transports specified power-sum orbit multiplicities to its
closure after a rectangular shift. Theorem4.3 gives the power-sum versus
product-of-variables multiplicity obstruction, with section5 treating
the stronger distinctions from occurrence tests. This transport theorem
does not directly identify the determinant orbit-closure restriction
image needed here.

\paragraph{Kumar and B\"urgisser--Ikenmeyer--H\"uttenhain.}
Kumar, \emph{A study of the representations supported by the orbit
closure of the determinant}, Compositio Mathematica151 (2015),292--312,
\href{https://kumar.math.unc.edu/papers/kumar67.pdf}{author PDF},
Theorem6.1 and Corollaries6.2,6.4,6.5, with Remark6.6(b).
The occurrence results include even-size/Alon--Tarsi hypotheses where
stated. B\"urgisser--Ikenmeyer--H\"uttenhain,
\emph{Permanent versus determinant: not via saturations},
\href{https://arxiv.org/abs/1501.05528}{arXiv:1501.05528},
Theorem1, treats saturation for partitions of length at most$n$ and
size divisible by$n$; Theorem2 states the conditional strengthening.
These statements do not supply A-030's length-$n^2$ explicit Hessian
equation or A-034's exact ray algebra. This is a scope comparison,
not an exhaustive novelty certificate for the inherited claims.

\paragraph{Landsberg's book.}
The inspected text is the author's 2016 prepublication version of
\emph{Geometry and Complexity Theory},
\href{https://people.tamu.edu/~jml/simonsclass.pdf}{author PDF};
its pagination differs from the published2017 edition.
Propositions8.4.1.1 and8.4.2.1--8.4.2.4 give subspace and padding
restrictions and explicitly use the independent-variable cap$m^2+1$.
Theorem8.5.2.1 and Theorems8.5.3.2--8.5.3.3 describe the LMR
dual-degeneracy equations and determinant component. Sections8.8--8.10
separate representation-theoretic approaches and occurrence limitations.
The book is cited, not redistributed in this package.

\paragraph{Beauville and determinantal surfaces.}
Beauville, \emph{Determinantal hypersurfaces}, Michigan Mathematical
Journal48 (2000),39--64, Theorem A, Proposition1.11 and Corollary1.12,
identify the ACM cokernel and its linear resolution; the converse in
Corollary1.12 concerns a smooth hypersurface and its line bundle.
These are the precise sheaf-theoretic inputs behind the proposed
Noether--Lefschetz route. Leal--Lozano Huerta--Vite,
\emph{Noether--Lefschetz locus and a special case of the variational
Hodge conjecture: using elementary surface theory},
\href{https://arxiv.org/abs/2303.09028}{arXiv:2303.09028v3},
Corollary3.1 gives the determinantal-locus dimension, and Theorem2/Table2
give degree320112 for the determinantal-quartic divisor. The latter
corrects the older doubled value640224; it is not a new computation in
this run and supplies no low-degree nonic defining equation.

\paragraph{Low-degree Chow equations.}
Cheung--Ikenmeyer--Mkrtchyan,
\href{https://arxiv.org/abs/1509.03944}{arXiv:1509.03944},
Theorem6(a) proves injectivity of the $(5,6)$ map; Theorem3 is the
McKay propagation result and Remark7 identifies the Chow ideal.
Together with lower-degree propagation these exclude Chow equations
through degree5 for$n\ge6$. Since diagonal determinants contain Chow,
this is background for starting a nonic determinant-equation search
in degree6. It is not a determinant-ideal calculation in degree6.

\paragraph{Border-rank audit in progress.}
The attached primary texts of Buczy\'nska--Buczy\'nski,
Conner--Huang--Landsberg, Landsberg--Micha\l ek and Farnsworth are
being checked against the arithmetic incidence proof. BR-001 specifies
the exact incidence statement used. A final priority claim for the
all-characteristic bound is withheld pending the complete targeted
literature register; Waring-rank and symmetric-border-rank bounds are
kept distinct from the CP tensor statement throughout.
